\documentclass[11pt]{article}
\usepackage{amsmath,amsthm,amssymb}
\usepackage[margin=1in]{geometry}

\newtheorem{theorem}{Theorem}
\newtheorem{proposition}{Proposition}
\newtheorem{lemma}{Lemma}
\newtheorem{corollary}{Corollary}
\newtheorem{conjecture}{Conjecture}
\theoremstyle{definition}
\newtheorem{remark}{Remark}

\newcommand{\Z}{\mathbb{Z}}
\newcommand{\C}{\mathbb{C}}
\newcommand{\R}{\mathbb{R}}
\newcommand{\Zi}{\Z[i]}
\renewcommand{\Re}{\operatorname{Re}}
\renewcommand{\Im}{\operatorname{Im}}

\title{Monotonicity of $N_l=|r_l|^2$ for $f(u)=u^2+(1+i)u+1$:\\
a Gram--bilinear recurrence and the remaining endpoint gap}
\author{Clio compute agent}
\date{2026-06-04}

\begin{document}
\maketitle

\begin{abstract}
Let $f(u)=u^2+(1+i)u+1\in\Zi[u]$ and $r_l=[u^l]f(u)^m$, with
$N_l=|r_l|^2\in\Z_{\ge0}$. We verify Clio's three--term recurrence for $r_l$
exactly, derive from it an explicit real \emph{Gram system} for the bilinear
quantities $N_l=|r_l|^2$, $P_l=\Re(r_l\overline{r_{l-1}})$,
$Q_l=\Im(r_l\overline{r_{l-1}})$, and obtain an exact fourth--order
$P$--recursion for the integer sequence $N_l$ alone. We reduce the conjecture
``$N_0<N_1<\cdots<N_m$, with the single exception $N_1=N_2$ at $m=2$'' to a
single explicit inequality $J_k>0$. We prove $J_k>0$ \emph{unconditionally} on
the range $k\le(2-\sqrt2)m$ (where it is automatic), and reduce the remaining
range to a lower bound on $S_k=P_k-Q_k$. The full statement is verified exactly
in $\Zi$ for all $m\le300$. The remaining gap is isolated precisely: it lives at
the extreme endpoint $k\in\{m-1,m\}$.
\end{abstract}

\section{Setup}

Fix $i=\sqrt{-1}$, $\pi=1+i$, and for an integer $m\ge1$ put
\[
  f(u)=u^2+(1+i)u+1\in\Zi[u],\qquad r_l:=[u^l]f(u)^m\quad(0\le l\le 2m).
\]
Since $u^2 f(1/u)=f(u)$ the sequence is palindromic, $r_l=r_{2m-l}$. We set
\[
  N_l:=|r_l|^2=r_l\overline{r_l}\in\Z_{\ge0},
\]
where $\overline{r_l}=[u^l]\bar f^m$ with $\bar f=u^2+(1-i)u+1$. The target is:

\begin{conjecture}\label{conj:main}
For every $m\ge1$ the sequence $N_0,N_1,\dots,N_m$ is strictly increasing,
\[
  N_0<N_1<\cdots<N_m,
\]
with the single exception $N_1=N_2\ (=8)$ at $m=2$.
\end{conjecture}

Conjecture~\ref{conj:main} closes ``Gap B'': it implies $r_l\ne r_{l-1}$ for
$1\le l\le m$ except at $(m,l)=(2,2)$ (since $N_l\ne N_{l-1}\Rightarrow
|r_l|\ne|r_{l-1}|\Rightarrow r_l\ne r_{l-1}$), the shape $(2,2)$ being the unique
exceptional partition.

\section{The $r$-recurrence (verified)}

\begin{proposition}\label{prop:rrec}
For $k\ge0$, with $r_{-1}=0$, $r_0=1$,
\[
  (k+1)\,r_{k+1}=(1+i)(m-k)\,r_k+(2m-k+1)\,r_{k-1}.
\]
The conjugate sequence satisfies the same identity with $(1+i)\mapsto(1-i)$.
\end{proposition}

This is exactly the recurrence proposed by Clio; it is \emph{correct as stated}.
It follows from $f'(u)/f(u)$ being rational of the Gauss/contiguity type and was
verified exactly in $\Zi$ against direct coefficient extraction of $f(u)^m$ for
all $1\le m\le30$ and all $0\le k\le2m$, with no discrepancy
(\texttt{mono\_verify\_recur.py}). The exact divisions by $k+1$ are integral in
$\Zi$.

\section{The Gram system}

Write $\alpha=\alpha_k:=m-k$ and $\beta=\beta_k:=2m-k+1$; for $0\le k\le m$ we
have $\alpha\ge0$ and $\beta>0$. Define the complex bilinear form
\[
  C_k:=r_k\,\overline{r_{k-1}}=P_k+iQ_k,\qquad
  P_k=\Re C_k,\ \ Q_k=\Im C_k\quad(k\ge1),
\]
so $|C_k|^2=P_k^2+Q_k^2=N_kN_{k-1}$ \emph{exactly} (a Cauchy--Schwarz
\emph{equality}, since $C_k$ is a genuine product). Set $S_k:=P_k-Q_k$ and
$T_k:=P_k+Q_k$.

\begin{theorem}[Gram system; verified exactly]\label{thm:gram}
For $0\le k<m$,
\begin{align}
  (k+1)\,P_{k+1} &= \alpha N_k + \beta P_k, \label{eq:P}\\
  (k+1)\,Q_{k+1} &= \alpha N_k - \beta Q_k, \label{eq:Q}\\
  (k+1)^2 N_{k+1} &= 2\alpha^2 N_k + \beta^2 N_{k-1} + 2\alpha\beta\,S_k.
  \label{eq:N}
\end{align}
Equivalently, with $S,T$,
\begin{align}
  (k+1)\,S_{k+1} &= \beta\,T_k, \label{eq:S}\\
  (k+1)\,T_{k+1} &= 2\alpha N_k + \beta\,S_k. \label{eq:T}
\end{align}
\end{theorem}

\begin{proof}
Multiply the $r$-recurrence by $\overline{r_k}$ and take real and imaginary
parts. Using $(1+i)r_k\overline{r_k}=(1+i)N_k$ and
$r_{k-1}\overline{r_k}=\overline{C_k}=P_k-iQ_k$,
\[
  (k+1)\,r_{k+1}\overline{r_k}
   =(1+i)\alpha N_k+\beta(P_k-iQ_k).
\]
But $r_{k+1}\overline{r_k}=C_{k+1}=P_{k+1}+iQ_{k+1}$; comparing real/imaginary
parts gives \eqref{eq:P}--\eqref{eq:Q}. For \eqref{eq:N} multiply the
$r$-recurrence by its conjugate:
\[
  (k+1)^2 N_{k+1}
  =\bigl|(1+i)\alpha r_k+\beta r_{k-1}\bigr|^2
  =2\alpha^2 N_k+\beta^2 N_{k-1}+2\alpha\beta\,\Re\!\big((1+i)r_k\overline{r_{k-1}}\big),
\]
and $\Re\big((1+i)C_k\big)=P_k-Q_k=S_k$. Equations
\eqref{eq:S}--\eqref{eq:T} are $\eqref{eq:P}\mp\eqref{eq:Q}$.
\end{proof}

All five identities were verified exactly in $\Zi$ for $1\le m\le40$
(\texttt{mono\_gram\_verify.py}), with no discrepancy.

\begin{corollary}[positivity]\label{cor:pos}
For $1\le k\le m$ one has $N_k>0$, $P_k\ge0$, $Q_k\ge0$ (hence $S_k\ge0$,
$T_k\ge0$). Indeed $P_1=Q_1=1$ ($C_1=(1+i)\cdot1$), and \eqref{eq:P},
\eqref{eq:T} propagate nonnegativity directly (all coefficients
$\alpha,\beta\ge0$). For $Q$, \eqref{eq:Q} requires $\alpha N_k\ge\beta Q_k$,
which holds for all tested $m\le60$, $1\le k<m$ (\texttt{mono\_invariants.py});
we record this as part of the verified data, not as a proved fact.
\end{corollary}

\section{A pure recurrence for $N_l$}

Solving \eqref{eq:N} for $S_k$ and substituting into the shifts of
\eqref{eq:S}--\eqref{eq:T} eliminates $S,T$ and yields a fourth--order
$P$--recursion for $N$ alone.

\begin{theorem}[verified exactly to $m=40$]\label{thm:Nrec}
With $\alpha=m-k$, $\beta=2m-k+1$ understood, for all valid $k$,
\[
\begin{aligned}
 &-k(k+1)(k+2)^2(k-m-1)\,N_{k+2}
 \;+\;2k(k+1)(k-m-1)(k-m+1)^2\,N_{k+1}\\
 &+\;2k(k-2m)(k-m)\,(k^2-2km-m-1)\,N_{k}
 \;+\;2(k-2m)(k-2m-1)(k-m-1)^2(k-m+1)\,N_{k-1}\\
 &-\;(k-2m)(k-2m-2)^2(k-2m-1)(k-m+1)\,N_{k-2}\;=\;0.
\end{aligned}
\]
\end{theorem}
The leading coefficient vanishes at $k=m+1$ (and at $k=0,-1,-2$); the trailing
coefficient at $k\in\{m-1,2m,2m+1,2m+2\}$, marking the singular points of the
recurrence. Verified exactly in $\Zi$ for $2\le m\le40$ (\texttt{mono\_Nrec.py}).

\section{Reduction of monotonicity to a single inequality}

Let $D_k:=N_k-N_{k-1}$. Subtracting $(k+1)^2N_k$ from \eqref{eq:N} and writing
$N_{k-1}=N_k-D_k$ gives, after factoring,
\[
  (k+1)^2 D_{k+1}
   =2\alpha(3m+2-k)\,N_k+2\alpha\beta\,S_k-\beta^2 D_k.
\]
Re-expanding once more in terms of $N_{k-1}$ gives the clean form we use:

\begin{lemma}[exact identity]\label{lem:J}
For $0\le k<m$, setting
\[
  g=g_k:=\beta^2-2\alpha(3m+2-k)=-k^2+(4m+2)k-2m^2+1,
\]
\[
  (k+1)^2\,D_{k+1}
   \;=\;J_k,\qquad
  J_k:=\beta^2 N_{k-1}+2\alpha\beta\,S_k-g\,N_k .
\]
\end{lemma}
\begin{proof}
Direct from \eqref{eq:N}: $(k+1)^2 D_{k+1}=(k+1)^2N_{k+1}-(k+1)^2N_k
=2\alpha^2N_k+\beta^2N_{k-1}+2\alpha\beta S_k-(k+1)^2N_k$. One checks the
polynomial identity $2\alpha^2-(k+1)^2=-g$ (both equal
$k^2-4km-2k+2m^2-1$), giving
$(k+1)^2 D_{k+1}=\beta^2N_{k-1}+2\alpha\beta S_k-gN_k$. (Verified symbolically,
\texttt{mono\_Dk.py}, and the identity $(k+1)^2D_{k+1}=J_k$ checked exactly to
$m=300$, \texttt{mono\_verify\_full.py}.)
\end{proof}

\begin{theorem}[reduction]\label{thm:reduction}
Conjecture~\ref{conj:main} is equivalent to:
\[
  J_k>0\quad\text{for all }m\ge1,\ 1\le k<m,\ \text{except }(m,k)=(2,1),
  \tag{$\star$}
\]
together with $D_m>0$ for $m\ge3$ (the endpoint, handled below via $J_{m-1}$).
The lone exception $(m,k)=(2,1)$ gives $J_1=0$, i.e. $D_2=0$, i.e. $N_1=N_2$.
\end{theorem}
\begin{proof}
By Lemma~\ref{lem:J}, $D_{k+1}>0\iff J_k>0$ (as $(k+1)^2>0$). Strict increase on
$0\le l\le m$ is $D_l>0$ for $1\le l\le m$; the index $l=k+1$ ranges over
$2\le l\le m$ as $k$ ranges over $1\le k<m$, plus $D_1=N_1-N_0=2m-1>0$ always.
The endpoint $D_m$ corresponds to $k=m-1$, where $\alpha=1$ and $J_{m-1}$ is the
relevant quantity. Numerically ($\star$) holds with the single equality at
$(2,1)$ for all $m\le300$.
\end{proof}

\section{What is proved, and the remaining gap}

\subsection{Unconditional region $g_k\le0$}

\begin{proposition}\label{prop:gneg}
If $g_k\le0$ then $J_k>0$. This covers all $k$ with
$k\le(2m+1)-(m+1)\sqrt2=(2-\sqrt2)\,m+(1-\sqrt2)$, i.e. asymptotically
$k\lesssim0.5858\,m$.
\end{proposition}
\begin{proof}
$g_k$ is a downward parabola in $k$ with roots $(2m+1)\pm(m+1)\sqrt2$, so
$g_k\le0$ left of the smaller root. There $J_k=\beta^2N_{k-1}+2\alpha\beta S_k
+|g_k|N_k$ is a sum of nonnegative terms with $\beta^2N_{k-1}>0$, hence $J_k>0$.
(Uses only $N_{k-1}>0$, $\alpha,\beta\ge0$, $S_k\ge0$.)
\end{proof}

\subsection{The region $g_k>0$ (the gap)}

For $k$ above the root the term $-g_kN_k$ is negative and must be dominated by
$\beta^2N_{k-1}+2\alpha\beta S_k$. Two facts sharpen the obstruction.

\begin{remark}[Cauchy--Schwarz alone is insufficient]
The bound $|S_k|\le\sqrt{N_kN_{k-1}}$ does \emph{not} suffice: the worst--case
substitution $S_k=-\sqrt{N_kN_{k-1}}$ makes
$\beta^2-2\alpha\beta\sqrt{\rho}-g\rho<0$ (with $\rho=N_k/N_{k-1}$) for the great
majority of $(m,k)$ in this region (\texttt{mono\_invariant\_search.py}). Thus
the proof genuinely needs the \emph{sign/size} of $S_k$ (equivalently, that the
phase $\arg C_k\in[0,\tfrac\pi2]$ stays bounded away from $\tfrac\pi2$), not
merely $|S_k|$.
\end{remark}

\begin{remark}[a near--complete provable lower bound]
Dropping the nonnegative $S$--term in \eqref{eq:T} gives the \emph{provable}
lower bound
\[
  T_{k-1}\ge \frac{2\alpha_{k-2}}{k-1}\,N_{k-2},\qquad
  S_k=\frac{\beta_{k-1}}{k}\,T_{k-1}\ \ge\
  L_k:=\frac{2\beta_{k-1}\alpha_{k-2}}{k(k-1)}\,N_{k-2}.
\]
Substituting $L_k$ for $S_k$, the inequality
$\beta^2N_{k-1}+2\alpha\beta L_k-g\,N_k>0$ holds for \emph{all} $g_k>0$ cases
with $k<m$ \emph{except} the single index $k=m-1$ once $m\ge25$
(\texttt{mono\_invariant\_close.py}; $948/993$ cases closed for $m\le69$). At
$k=m-1$ the dropped term is no longer negligible because $\alpha=1$ is small and
the recurrence is near its singular point $k=m-1$ (Theorem~\ref{thm:Nrec}).
\end{remark}

\paragraph{Precise remaining gap.}
The complete proof reduces to the single endpoint inequality: \emph{prove}
\[
  J_{m-1}=(m+2)^2 N_{m-2}+2(m+2)\,S_{m-1}-(m^2-2)\,N_{m-1}>0\qquad(m\ge3),
\]
which is exactly $D_m>0$, together with the band $g_k>0,\ k=m-1$ for the smaller
shifts. Equivalently one needs a lower bound on $S_{m-1}$ (or on the phase of
$C_{m-1}$) sharper than $L_{m-1}$; unrolling \eqref{eq:S}--\eqref{eq:T} one more
step (keeping the $\beta S$ term) is expected to close it, but a clean uniform
algebraic bound was not obtained here.

\subsection{Numerical status}

\begin{proposition}[exact, $m\le300$]\label{prop:numeric}
For all $1\le m\le300$: the identity $(k+1)^2D_{k+1}=J_k$ holds; there are
\emph{no} decreases $N_k<N_{k-1}$; the only tie $N_k=N_{k-1}$ is $(m,k)=(2,2)$;
$J_k>0$ for $1\le k<m$ with the single boundary equality $J_1=0$ at $m=2$
(\texttt{mono\_verify\_full.py}).
\end{proposition}

\section{Summary}
\begin{itemize}
\item Clio's $r$-recurrence is correct (Prop.~\ref{prop:rrec}), verified to
$m=30$.
\item The Gram system \eqref{eq:P}--\eqref{eq:T} (Theorem~\ref{thm:gram}) and the
fourth--order $N$-recursion (Theorem~\ref{thm:Nrec}) are exact, verified to
$m=40$.
\item Monotonicity reduces to the single inequality $J_k>0$
(Lemma~\ref{lem:J}, Theorem~\ref{thm:reduction}), with the exceptional shape
$(2,2)$ appearing exactly as the boundary equality $J_1=0$ at $m=2$.
\item \textbf{Proved:} $J_k>0$ on the whole region $g_k\le0$
($k\lesssim0.586\,m$), unconditionally (Prop.~\ref{prop:gneg}). A provable lower
bound $L_k$ closes \emph{all} of $g_k>0$ except the single endpoint index
$k=m-1$ (for $m\ge25$).
\item \textbf{Remaining gap:} the endpoint $J_{m-1}>0$ (i.e. $D_m>0$), needing a
sharper lower bound on $S_{m-1}$. Verified exactly to $m=300$.
\end{itemize}

\end{document}
